\section{Theorie}\label{sec:theorie}
\subsection{Adsorbtionsmittel zur Anreicherung der Analyten}
Voraussetzung für die Quantifizierbarkeit der emittierten Verbindungen ist, diese in erfassbaren Konzentrationsbereichen so verlustfrei wie möglich anzureichern.
Zur Verfügung stehen dafür die physikalische und chemische Adsorbtion in Flüssigkeiten, die Kondensation bei tiefen Temperaturen, sowie die Adsorbtion an Festkörpern.
Letztere hat sich für die Spurenanalytik von organischen Luftinhaltsstoffen aufgrund des heutzutage bestehenden großen Angebots an festen Sorbensmaterialien etabliert. \par
Genannte feste Sorbenzien lassen sich in anorganische Adsorbtionsmittel, Adsorbtionsmittel auf Kohlenstoffbasis und poröse Polymere einordnen.
Unerlässliche Anforderungen sind dabei, dass das gewählte Adsorbtionsmittel bei Umgebungstemperatur ein großes \ac{vd} für organische Verbindungen wie z.B. \gls{btex} aufweist, die Analyten nach der Adsorbtion wieder vollständig zu eluieren vermag und während beider Prozessbedingungen eine hohe physikalische und chemische Stabilität aufweist. \par
Das \ac{vd}, angegeben in Liter Gas pro Gramm Adsorbens, ist dabei eine auf die Adsorbensmasse bezogene Größe, welche die Gasmenge angibt die das Adsorbtionsröhrchen durchfließen darf bevor eine Spezies eluiert. 
\ac{vd} ist stark temperaturabhängig und individuell für jede adsorbierte Spezies \cite{figge_adsorptionsmittel_1987}. \par
Bei der Wahl des Adsorbens kommt es dabei vor allem auf die Flüchtigkeit des Analyten an.
So sind für sehr flüchtige Verbindungen starke, für halbflüchtige Verbindungen eher schwache Adsorbens Materialien geeignet.
